\section{Construction of the complexity reasoning graph collection}\label{app:complexity-graphs}

This appendix describes the complexity atlas through October 31, 2026, including the later exploration component from which the certificate-degree investigation drew its inputs. The October 25 extraction is the base of this evolving collection. It separates three objects: the extracted \emph{complexity atlas}, the selected \emph{working graph} used for the certificate-degree investigation, and the \emph{increment} produced by that investigation. The atlas implements source reconstruction; the run illustrates research that uses and extends selected reasoning. Neither the collection size nor its extraction status measures the amount of independently verified mathematics.

\subsection{Sources, snapshots, and units of accounting}

The atlas is a local, dated collection of original complexity research notes, manuscript sections, arguments, research histories, static code contracts, and stored computational results. Its October 25, 2026 completion record covers five local research collections spanning complexity notes, problem directions, prior exploratory records, and focused investigations. Source location records provenance, not authorship. Separate research streams retain their original files and hypotheses.

\Cref{tab:atlas-evolution} distinguishes the original extraction from the atlas available when the present case study began and the atlas after its final integration. Between the extraction snapshot and the new run, subsequent components added 4,578 node records and 2,667 AND inferences. The case study then contributed a component of 68 nodes and 34 AND inferences.

\begin{table}[htbp]
\centering\small
\begin{tabularx}{\linewidth}{@{}Yrr@{}}
\toprule
Dated atlas state & Nodes & AND inferences \\
\midrule
October 25: original extraction & 58,916 & 19,516 \\
October 31: before the case study & 63,494 & 22,183 \\
October 31: after the case study & 63,562 & 22,217 \\
\bottomrule
\end{tabularx}
\caption{Cumulative atlas counts through October 31, 2026. Each row is a complete dated state, not an additional component. Only the registered version of each component is counted, and references to existing base nodes are not counted again.}
\label{tab:atlas-evolution}
\end{table}

The document and batch counts below describe the original extraction; later research components are accounted for by their graph records. The extraction denominator excludes 49 adjacent collections, downloaded libraries, Git internals, and generated infrastructure. External theorem statements may appear as imports; their full proofs are not thereby extracted. Completion means that the declared snapshot's in-scope files received the recorded extraction or documentary treatment, with no pending or stale in-scope files in that audit. It does not describe later additions or every research subject in the workspace.

\begin{table}[htbp]
\centering\small
\begin{tabularx}{\linewidth}{@{}Yrp{5.8cm}@{}}
\toprule
Unit & Count & Interpretation \\
\midrule
Primary documents & 1,874 & Source files or sections, not necessarily papers \\
Reconciled history documents & 855 & Earlier records related to later argument records \\
Extraction batches & 4,200 & Units of annotation and accounting \\
Mathematical argument groups & 7,738 & Local groups of claims, scopes, and steps \\
History and support groups & 17,406 & Documentary, review, and supporting records \\
Nodes & 58,916 & Multiple semantic roles, not only theorems \\
AND inferences & 19,516 & Authored joint-premise derivations \\
Cross-batch navigation links & 15,515 & Connections that do not assert entailment \\
Distinct source-file pins & 12,743 & Primary and supporting file references \\
\bottomrule
\end{tabularx}
\caption{The complexity atlas snapshot completed on October 25, 2026. Counts are different units and should not be added as if they measured independent results.}
\label{tab:atlas-counts}
\end{table}

The 4,200 batches include 1,691 with local inferences and 2,509 without them. Definitions, imported contracts, historical records, and file-accounting groups can legitimately have no local derivation. Of the nodes, 19,308 have role \code{derived_claim}, 18,207 \code{review_record}, 967 \code{open_obligation}, and 244 \code{counterexample}. These roles describe recorded content. In particular, a source-reported derived claim is not automatically a theorem accepted by RRE.

\subsection{Semantic extraction and assembly}

The annotations were drafted by the AI assistant through reading and explicit semantic authoring. The compiler assembles those annotations and checks their structure; it does not autonomously reconstruct arguments from unannotated prose. The procedure is:

\begin{enumerate}[leftmargin=1.7em]
\item \textbf{Inventory and pin the sources.} Record file paths, SHA-256 hashes, source roles, and immutable snapshots. Separate mathematical documents from histories, supporting programs, saved results, and excluded infrastructure.
\item \textbf{Identify argument units.} Select a conclusion and reconstruct its supporting steps from the source. Paragraph order, common vocabulary, and citation links alone do not establish a dependency.
\item \textbf{Expose the hypotheses.} Record the field, domain, quantifier order, fixed and growing parameters, degree and support bounds, exceptional cases, and distinctions such as exact versus border statements. A proposed bridge remains conditional when its hypotheses have not been shown for the intended circuit or target.
\item \textbf{Author claims and AND steps.} Give each claim a role, source passage, and scope. List jointly required premises explicitly. Keep definitions, external contracts, computational assertions, counterexamples, and unresolved obligations distinct.
\item \textbf{Reconcile and connect.} Relate historical attempts to later arguments without upgrading old confidence statements into proofs. Add cross-batch navigation for relevant tools, obstructions, and possible transfers; retain the proof obligations for using those connections.
\item \textbf{Compile and check.} Assemble batch graphs, propagate recorded scope, check identifiers and evidence anchors, test prerequisite traces, account for source coverage, and export the navigable collection.
\end{enumerate}

Source programs were read statically during this extraction. A saved exact computation retains the status of a computation reported by its source unless a separately identified replay was performed. The fresh checks of the certificate-degree investigation in \cref{sec:case} are a different operation from extracting the atlas.

\subsection{Implemented records and scope propagation}

The implemented atlas schema is narrower than the proposed RRE record specification in \cref{app:records}. Its principal fields are summarized below; the companion sample includes the actual JSON schema and representative records.

\begin{table}[htbp]
\centering\small
\begin{tabularx}{\linewidth}{@{}p{2.2cm}Y@{}}
\toprule
Record & Implemented information \\
\midrule
Source & Path, full-file SHA-256, snapshot path, source role \\
Evidence anchor & Source ID, inclusive line interval, exact decoded passage, passage SHA-256 \\
Node & Stable ID, authored title, verbatim source passages, role, source-reported status, scope IDs, review fields, planner eligibility \\
Inference & ID, premise IDs, conclusion, scope IDs, effective premises, reasoning mode, evidence; explicit assumption discharge where needed \\
Argument group & Nodes, targets, inferences, scope, and interpretation of conditional hypotheses \\
Navigation relation & Endpoints, relationship kind, explanatory note, and \code{asserts_entailment=false} \\
\bottomrule
\end{tabularx}
\caption{Existing atlas records. An authored summary and its verbatim evidence remain separate.}
\end{table}

For an inference, the compiler begins with its explicit premises, its own scope, and the conclusion's attached scope. It then repeatedly adds scopes attached to these nodes to produce \code{effective_premises}. This prevents a recorded hypothesis from disappearing merely because a viewer displays only the immediate premise list. It cannot recover a hypothesis that the annotation omitted. When an argument discharges a local assumption, that event is recorded separately; the assumption's appearance in the historical prerequisite trace does not make it an undischarged hypothesis of the conclusion.

All nodes and inference records in the original extraction have \code{planner_eligible=false}. Their review fields distinguish source grounding from mathematical certification and say that human review was not assessed in this import. RRE must therefore select, interpret, and assess a proposed input before treating it as supported, or retain it explicitly as a conditional source dependency. Importing a graph does not silently admit all its assertions into a proof engine.

\subsection{A concrete conditional argument}

Batch CB4173 illustrates why this representation matters. Its source is a dated transfer-and-proofs note in the atlas integration materials, lines 129--147. The scope fixes source coefficient polynomials $Q_i$ with identically zero off-support coefficients, target polynomials $\mu_i$, and the \emph{full} space of coefficient monomials of degree at most $D$. Let
\[
 A_D(P)=P(Q),\qquad T_D(P)=P(\mu).
\]
The node \code{CB4173:product-blind} is the hypothesis $\ker A_D\subseteq\ker T_D$. Under that hypothesis,
\[
 \Phi_D(P(Q))=P(\mu),\qquad \deg P\le D,
\]
is well defined and linear: two representatives with the same source substitution differ by an element of $\ker A_D$, which the hypothesis places in $\ker T_D$. If $\deg P\le a$, $\deg R\le b$, and $a+b\le D$, substitution gives
\[
 \Phi_D\bigl(P(Q)R(Q)\bigr)
   =\Phi_D(P(Q))\Phi_D(R(Q)).
\]
The product is available because the full degree-bounded space contains $PR$.

The recorded step \code{CB4173:multiplicative-infer} has the three premises \code{product-scope}, \code{product-blind}, and \code{truncated-map}, all in the CB4173 namespace. Thus a trace exposes the full-dictionary scope and the kernel-inclusion hypothesis. It does not present multiplication compatibility as an unconditional consequence of a sparse linear reconstruction. Separate nodes retain the cost of the full dictionary and the unresolved obligation to obtain an affordable extension for actual circuit sources. This is an example of an existing conditional argument, not an additional discovery of the present investigation. The short derivation and its source anchors were checked during preparation of this appendix; no broader novelty or human-review claim follows.

\subsection{From the atlas to an evolving working graph}

The case study starts from the October 31 atlas, not solely from the October 25 extraction. Its immediate source is revision 241 of the preceding exploration component, RRE003, with 4,108 nodes and 2,362 AND inferences. The new run pins this graph and its supporting note, then selects two results: the distinct-label lower functional and the matching certificate when a flattening has rank greater than the proposed tensor rank. Their canonical records are \code{RRE003:cp119-lower} and \code{RRE003:cp119-exact}. These become explicit inherited inputs in the new component, RRE005; their hypotheses are checked for the new family. Thus the case study reuses results of an earlier exploration already incorporated into the atlas. Earlier investigations of coefficient reconstruction and parameter normalization, related to scaling reparametrizations such as \cite{meshkat}, also remain available in the collection.

The new component starts with eight records identifying its objective, parameter convention, inherited results, and boundary target. Successive publications add proofs, finite certificates, generated targets, and classical imports. Its six revisions contain respectively 8, 19, 43, 46, 66, and 68 nodes, and 0, 7, 22, 22, 33, and 34 AND inferences. The final component includes its eight initial records, so growth after initialization is 60 nodes and 34 inferences. Each publication stores the complete component, a delta from its predecessor, and validation results, then advances only this component's atlas pointer. Other components retain their own versions.

At the final integration on October 31, the registry reaches the last row of \cref{tab:atlas-evolution}, with 17,663 navigation links across the base and subsequent components. This navigation total includes both local and cross-component relations; the 15,515 links in \cref{tab:atlas-counts} count only the original cross-batch relations. The 68-node component is the object analyzed in the case study, while the larger atlas supplies the research context and selected inherited inputs. Two dedicated discovery exports retain the full prerequisite paths for the nilpotent-pencil theorem and the exact two-term theorem, with copies of their hash-pinned evidence. This permits a separate visualization process to inspect a result without copying the entire atlas.

The collection viewer serves corpus navigation; the certificate-study companion serves the selected investigation. The latter offers a saved view for each version, complete recorded prerequisite traces, AND-step inspection, and source evidence. The printed overview in \cref{fig:certificate-versions} relates the beginning and end of this component. Its timeline in \cref{fig:discovery-timeline} describes research actions, while the theorem view in \cref{fig:certificate-proof} describes implication.

\subsection{Validation, limitations, and companion material}

The dated atlas audit reports 105,294 evidence anchors checked, 2,725 prerequisite regressions, and automated browser checks of 25,144 argument groups and their expanded traces. These checks concern source alignment, expected graph behavior, and interface operation. Nonblank-line coverage also includes documentary classifications and broad source spans; it is not a measure of semantic accuracy. Neither automated coverage nor a clean graph certifies every inference or establishes novelty.

For the original extraction appendix, the full graph's SHA-256 was checked against the completion and statistics reports. Batch totals, node roles and statuses, anchor counts, and planner-ineligibility flags were recounted from all 4,200 batch graphs. The sources, snapshots, and evidence anchors of CB291 and CB4173 were checked again. These read-only checks did not rerun all historical regressions, browser tests, or source computations. The new component has separate source-hash, dependency, exact-certificate, and viewer checks, documented in \cref{app:records}.

For the October 31 accounting, the base graph and all five registered components were checked against their recorded SHA-256 hashes. Node and inference counts were recomputed, repeated base references were excluded, and the totals were matched to the registry saved before the case study and the registry after its final integration. The two direct imports were traced to the pinned prior graph and proof note. The companion accounting report retains these source pins, the two registry snapshots, and the recount program.

The local atlas is organized as a compiled semantic graph, with batch graphs and authored specifications stored separately. The paper package supplies a representative atlas sample containing the extraction documentation, schema, dated reports, and two batches with their sources. The accompanying source manifest records paths and SHA-256 hashes, and the audit report separates newly performed checks from historical reported checks. The full atlas is not bundled with this draft. This sample documents the method and its limits without implying that the entire collection has been publicly released.
